% TOP_PROBLEM: {"id": 3213, "title": "Good Edge Labelings", "category_id": 3, "category": {"id": 3, "name": "graph_theory", "display_name": "Graph Theory", "description": "Problems involving graphs, networks, and their properties.", "slug": "graph-theory", "order_index": 3, "created_at": "2026-07-31T15:26:25.671Z"}, "status": "open", "problem_number": "OPG-37323", "set_id": 12}
% FIRST_POSTED: 2026-10-01
% ENTRY_KIND: statement_only
% UnsolvedMath ID 3213; source catalogue code OPG-37323
% !TeX program = lualatex
% Definition and sources only; this file is not an LLM solution attempt.
\documentclass[11pt,a4paper]{article}
\usepackage[margin=25mm]{geometry}
\usepackage{fontspec}
\setmainfont{lmroman10-regular.otf}[BoldFont=lmroman10-bold.otf,
 ItalicFont=lmroman10-italic.otf,BoldItalicFont=lmroman10-bolditalic.otf]
\usepackage{amsmath,amssymb,mathtools}
\usepackage{unicode-math}
\setmathfont{latinmodern-math.otf}
\AtBeginDocument{\let\setminus\smallsetminus}
\usepackage{xurl,hyperref}
\hypersetup{colorlinks=true,urlcolor=blue}
\setcounter{secnumdepth}{0}
\setlength{\parindent}{0pt}
\setlength{\parskip}{5pt}
\setlength{\emergencystretch}{3em}
\begin{document}
\section*{Good Edge Labelings}\label{3213}
{\small UnsolvedMath ID 3213 · OPG-37323 · Graph Theory · OpenGarden}
\subsection{Definitions and mathematical statement}
Let $G$ be a finite simple graph. An edge labeling is an injective map $\lambda:E(G)\to\mathbb R$. On a cycle, an edge is a local maximum if its label exceeds the labels of both neighboring edges of that cycle. Call the labeling good if every cycle has at least two local maxima, and call the graph good if it has some good labeling. Equivalently, between any ordered pair of vertices there is at most one path whose successive edge labels strictly increase.

The first problem is to determine the extremal function
\[
 f(n)=\max\{|E(G)|:|V(G)|=n\text{ and }G\text{ has a good labeling}\},
 \qquad n\ge1.
\]
The choice of real labels is unconstrained apart from being distinct. Only their relative order affects goodness, so one may equivalently use the integers $1,\ldots,|E(G)|$ in some order.

Call a graph critical if it is not good but every proper subgraph is good. Here a proper subgraph may be obtained by deleting vertices, edges, or both; criticality is not restricted to induced subgraphs. The second question asks whether, for every real constant $c<4$, there are only finitely many critical graphs, up to isomorphism, with average degree $2|E(G)|/|V(G)|<c$.

These are different targets: one is a maximum-density problem among good graphs, while the other concerns minimal obstructions below each fixed density threshold. A bound depending on $c$ is allowed in the finiteness assertion.
\subsection{Short English statement}
Determine the maximum size of a graph with a good edge labeling, and whether only finitely many minimal bad graphs occur below any fixed average degree less than four.
\subsection{Sources}
\begin{itemize}
\item[S1] ULAM.AI and the UnsolvedMath Contributors, supplied problems.json, record 3213 (OPG-37323), OpenGarden. Catalogue identity and source transcription; CC BY 4.0.
\url{https://huggingface.co/datasets/ulamai/UnsolvedMath}.
\item[S2] Open Problem Garden, Good Edge Labelings. Original problem and discussion.
\url{http://www.openproblemgarden.org/op/good_edge_labelings}.
\item[S3] A. Mehrabian, D. Mitsche and P. Prałat, On the Maximum Density of Graphs with Good Edge-Labellings, SIAM J. Discrete Math. 27 (2013), 1228--1233.
\url{https://arxiv.org/abs/1211.2641}.
\end{itemize}
\end{document}
